\begin{longtable}[c]{|l|m{10cm}|}
    \caption{运算指令}
    \label{tab:pie-vmul-instructions}\\
    \hline
    \rowcolor{lightgray}
    \textbf{指令} & \textbf{定义} \\\hline
    \endfirsthead
    \rowcolor{lightgray}
    \textbf{指令} & \textbf{定义} \\\hline
    \endhead
    \endfoot

    \hyperref[instruction:ESPVMULU8]{ESP.VMUL.U[8/16]} & 对无符号1字节/2字节数据执行向量乘法操作。 \\\hline
    \hyperref[instruction:ESPVMULS8]{ESP.VMUL.S[8/16]} & 对有符号1字节/2字节数据执行向量乘法操作。 \\\hline
    \hyperref[instruction:ESPVMULU8LDINCP]{ESP.VMUL.U[8/16].LD.INCP} & 对无符号1字节/2字节数据执行向量乘法操作，同时从内存读取16字节数据。 \\\hline
    \hyperref[instruction:ESPVMULS8LDINCP]{ESP.VMUL.S[8/16].LD.INCP} & 对有符号1字节/2字节数据执行向量乘法操作，同时从内存读取16字节数据。 \\\hline
    \hyperref[instruction:ESPVMULU8STINCP]{ESP.VMUL.U[8/16].ST.INCP} & 对无符号1字节/2字节数据执行向量乘法操作，同时将16字节数据写入内存。 \\\hline
    \hyperref[instruction:ESPVMULS8STINCP]{ESP.VMUL.S[8/16].ST.INCP} & 对有符号1字节/2字节数据执行向量乘法操作，同时将16字节数据写入内存。 \\\hline
    \hyperref[instruction:ESPVMULS32S16XS16]{ESP.VMUL.S32.S16xS16} & 执行16位向量乘法操作并右移结果。
    \\\hline
\end{longtable}